\subsection*{Metamathematics}

Absence from contradiction has always been considered to be a sine qua non of all mathematics, and from the time of Aristotle logic was sufficiently developed for it to be realized that anything could be deduced from a contradictory theory. Proofs of ``existence'', which have been regarded as indispensable ever since antiquity, clearly served only to guarantee that the introduction of a new concept did not risk entailing a contradiction, particularly if the concept was too complicated to fall immediately under ``intuition''. We have seen how this demand for a proof of non-contradiction became more imperative with the advent of the axiomatic point of view in the nineteenth century, and how the construction of arithmetic ``models'' answered it. But might arithmetic itself not be contradictory? Such a question could not have been asked before the end of the nineteenth century; so much did the integers appear to belong to what is surest in our intuition. But after the ``paradoxes'' everything seemed open to doubt, and the feeling of insecurity they created understandably led mathematicians, round about 1900, to look more closely at the problem of the consistency of arithmetic, in order to save at least classical mathematics from shipwreck. Thus this problem was the second of those listed by Hilbert in his famous address to the International Congress of 1900 ({[}31{]}, pp.~229-301). In doing so, he put forward a new principle which was to make a deep impression. Whereas, in traditional logic, the non-contradiction of a concept only made the concept ``possible'', for Hilbert it was equivalent to the existence of the concept (at any rate where axiomatically defined mathematical concepts were concerned). This point of view apparently entailed the necessity of proving a priori the consistency of a mathematical theory before even being able to develop the theory legitimately; certainly this was how the principle was understood by H. Poincaré, who took up Hilbert's idea as a convenient stick with which to beat the formalists, by pointing out with malicious satisfaction how far the formalists in this period were from being able to satisfy this condition ({[}33 c{]}, p.~163). We shall see later how Hilbert took up this challenge. But first, it should be noted that, under his influence and that of Poincaré, the demands put forward by the latter were accepted without reserve for a long time, as much by the formalists as by their opponents. One consequence was the belief, which was very widespread even among the formalists, that Hilbert's theory of proof was an integral part of mathematics, and constituted an indispensable prolegomenon to mathematics proper. We explained in the Introduction why this dogma appears unjustified to us (*), and we hold that the role of metamathematics in an account of logic and mathematics can and should be confined to the very elementary part which deals with the manipulation of abbreviating symbols and deductive criteria. Contrary to what Poincaré asserted, it is thus not a question of ``claiming freedom from contradiction'', but rather of considering, as Hadamard did, that consistency, even when it cannot be proved, can be ascertained ({[}35{]}, p.~270).

It remains for us to give a brief historical sketch of the efforts of Hilbert and his school. Although the theory of proof is not touched on in this series, it is not without interest to retrace in outline not only the evolution of ideas which finally led to Gödel's negative result and justified a posteriori

Hadamard's scepticism, but also all the progress which has resulted in an understanding of the mechanism of mathematical reasoning, and which has elevated modern metamathematics to the position of an autonomous science of undeniable importance.

In 1904, in an address to the International Congress ({[}30{]}, pp.~247-261), Hilbert attacked the problem of the consistency of arithmetic. First, he established that a proof of consistency could not be obtained by recourse to a ``model'' (*), and he indicated in broad outline the principle of another method. He proposed to consider the true propositions of formalized arithmetic as assemblies of signs without any meaning and to show that, by using the rules governing the formation and juxtaposition of these assemblies, an assembly could never be obtained which was a true proposition and whose negation was also a true proposition. He sketched a proof of this nature for a formalism less elaborate than that of arithmetic; but, as H. Poincaré observed soon afterward ({[}33 c{]}, p.~185), this proof made essential use of the principle of induction, and therefore appeared to be founded on a vicious circle. Hilbert did not reply to this criticism immediately, and fifteen years went by before anyone tried to develop his ideas. It was only in 1917 that (moved by the desire to answer the attacks of the intuitionists) he devoted himself again to the problem of the foundations of mathematics, which from then on occupied him continuously until the end of his scientific career. In his work on this subject, which stretches from 1920 to about 1930, and in which a whole school of young mathematicians (Ackermann, Bernays, Herbrand, von Neumann) were active participants, Hilbert gradually extracted the principles of his ``theory of proof'' in a more precise fashion. Realizing implicitly the justice of Poincaré's criticism, he admitted that in metamathematics the arithmetical arguments used could be based only on our intuition of the integers (and not on formalized arithmetic). Thus it appeared essential to restrict these arguments to ``finite procedures''(finite Prozesse) of a type allowed by the intuitionists. For example, a proof by reductio ad absurdum cannot establish the metamathematical existence of an assembly or a sequence of assemblies; it is necessary to give an explicit law of construction (†). Furthermore, Hilbert enlarged his initial program in two directions; not only did he attack the problem of the consistency of arithmetic, but he also hoped to prove the consistency of the theory of real numbers and even of the theory of sets (*). To the problems of consistency were added those of independence of the axioms, of categoricity, and the decision problem. We shall give a brief review of these various questions and mention the principal researches which they have occasioned.

A proof of the independence of a system of propositions \(A_{1}, A_{2}, \ldots, A_{n}\) consists in showing that, for each index i, \(A_{i}\) is not a theorem in the theory \(G_{i}\) obtained by taking as axioms the \(A_{j}\) where \(j \neq i\) . All we need is to find a non-contradictory theory \(G_{i}^{\prime}\) in which the \(A_{j} (j \neq i)\) are theorems, together with ``not \(A_{i}\) ''; we may therefore consider the problem in two aspects according as we do or do not assume that certain theories (such as arithmetic or set theory) are consistent. In the latter case, we have a problem of ``absolute'' consistency. On the other hand, the first type of problem is to be solved, like problems of ``relative'' consistency, by the construction of appropriate ``models'', and many proofs of this nature were devised well before mathematics had taken on a completely formalized aspect. We may cite as examples the models of non-Euclidean geometry, the questions of the independence of the axioms of elementary geometry treated by Hilbert in the Grundlagen der Geometrie {[}30{]}, the work of Steinitz on the axiomatization of algebra, and the work of Hausdorff and his successors on the axiomatization of topology.

A theory \(\mathfrak{C}\) is said to be categorical if, for every proposition A in \(\mathfrak{C}\) which contains no letters other than the constants of \(\mathfrak{C}\) , one of the two propositions A, (not A) is a theorem in \(\mathfrak{C}\) ( \(\dagger\) ). Apart from some very rudimentary formalism whose categoricity is easily proved ({[}32{]}, p.~35; cf.~Chapter I, Appendix, Exercise 7), the results obtained in this area are essentially negative. The most important is due to K. Gödel who showed that if \(\mathfrak{C}\) is non-contradictory and if the axioms of formalized arithmetic are theorems in \(\mathfrak{C}\) , then \(\mathfrak{C}\) is not categorical. The fundamental idea of his ingenious method consists in establishing a one-to-one correspondence (of course, by means of ``finite procedures'') between metamathematical statements and certain propositions of formalized arithmetic. We give here an outline sketch of the argument (**). With each assembly A which is a term or a relation in \(\mathfrak{C}\) , we associate (by an explicit constructive procedure which can be applied quasi-mechanically) an integer \(g(A)\) in a one-to-one manner. Likewise, with every proof D in \(\mathfrak{C}\) (considered as a succession of assemblies; cf.~Chapter I, § 2, no. 2) we may associate an integer \(h(\mathbf{D})\) in a one-to-one manner. Finally, we can give an explicit procedure for constructing a relation \(\mathrm{P}\{x, y, z\}\) in \(\mathcal{C} (*)\) such that, in \(\mathcal{C}, \mathrm{P}\{x, y, z\}\) implies that \(x, y, z\) are integers, and satisfies the following two conditions:

\begin{enumerate}
\def\labelenumi{(\arabic{enumi})}
\tightlist
\item
  If \(\mathbf{D}\) is a proof of \(\mathbf{A}\{\lambda\}\) , where \(\mathbf{A}\{x\}\) is a relation in \(\mathfrak{G}\) , and \(\lambda\) is a definite integer (that is, a term in \(\mathfrak{G}\) which is an integer), then
\end{enumerate}

\[
\mathbf {P} \{\lambda , g (\mathbf {A} \{x \}), h (\mathbf {D}) \}
\]

is a theorem in \(\mathfrak{G}\) .

\begin{enumerate}
\def\labelenumi{(\arabic{enumi})}
\setcounter{enumi}{1}
\tightlist
\item
  If the definite integer \(\mu\) is not of the form \(h(\mathbf{D})\) , or if \(\mu = h(\mathbf{D})\) and \(\mathbf{D}\) is not a proof of \(\mathbf{A}\{\lambda\}\) , then (not \(\mathbf{P}\{\lambda, g(\mathbf{A}\{x\}), \mu\}\) ) is a theorem in \(\mathfrak{C}\) .
\end{enumerate}

Now let \(S\{x\}\) be the relation (not \((\exists z)P\{x, x, z\}\) ), and let \(\gamma = g(S\{x\})\) , which is a term in \(\mathfrak{C}\) . If \(\mathfrak{C}\) is not contradictory, there is no proof of the proposition \(S\{\gamma\}\) in \(\mathfrak{C}\) . For if D were such a proof, then

\[
\mathrm{P} \{\gamma , g (\mathrm{S} \{x \}), h (\mathrm{D}) \}
\]

would be a theorem in \(\mathcal{G}\) . This relation is just \(\mathrm{P}\{\gamma, \gamma, h(\mathrm{D})\}\) , and consequently \((\exists z)\mathrm{P}\{\gamma, \gamma, z\}\) would also be a theorem in \(\mathcal{G}\) ; but since the last relation is equivalent to (not \(\mathrm{S}\{\gamma\}\) ), \(\mathcal{G}\) would be contradictory. Furthermore, what has just been said shows that, for each definite integer \(\mu\) , (not \(\mathrm{P}\{\gamma, \gamma, \mu\}\) ) is a theorem in \(\mathcal{G}\) . It follows that there is no proof in \(\mathcal{G}\) of the proposition (not \(\mathrm{S}\{\gamma\}\) ), for this relation is equivalent to \((\exists z)\mathrm{P}\{\gamma, \gamma, z\}\) , and the existence of an integer \(\mu\) such that \(\mathrm{P}\{\gamma, \gamma, \mu\}\) would imply that \(\mathcal{G}\) was contradictory, by virtue of the preceding argument ( \(\dagger\) ). This metamathematical theorem of Gödel has been subsequently generalized in various directions ({[}48{]}, Chapter XI) (**).

The relation \(S\{\gamma\}\) in \(\mathcal{C}\) , which is thus shown to have the property that there exists no proof in \(\mathcal{C}\) either of it or of its negation, has obviously been manufactured to fit the requirements of the argument and is not related in any natural way to any mathematical problem. What is much more interesting is the fact that, if \(\mathcal{C}\) denotes the theory of sets (with the system of axioms of von Neumann-Bernays), neither the continuum hypothesis nor its negation are provable in \(\mathcal{C}\) . This remarkable result has been established in two steps: in 1940, Gödel proved that the theory obtained by adjoining to \(\mathcal{C}\) the hypothesis \(2^{\aleph_0} = \aleph_1\) is not contradictory {[}44b{]}; and quite recently P. Cohen has proved that the same is met when the relation \(2^{\aleph_0} = \aleph_2\) (or \(2^{\aleph_0} = \aleph_n\) for any integer \(n > 1\) ) is adjoined to \(\mathcal{C}\) {[}49{]}.

The decision problem (Entscheidungsproblem) is certainly the most ambitious of all in metamathematics. The problem is whether, for a given formalized language, a quasi-mechanical ``universal procedure'' can be conceived which, when applied to any relation whatever in the formalism under consideration, will indicate in a finite number of operations whether the relation is true or not. The solution of this problem formed a substantial part of the grand designs of Leibniz, and it seems that at one point the Hilbert school believed they were very close to a solution. It is a fact that such procedures can be described for formalisms which contain few primitive signs and axioms ({[}48{]}, pp.~136-141; cf.~Chapter I, Appendix, Exercise 7). But the efforts to make precise the decision problem by defining exactly what is to be meant by ``universal procedure'' have led so far only to negative results ({[}48{]}, pp.~432-439). Moreover, the solution of the decision problem for a theory G would tell immediately whether or not G is contradictory, because the ``universal procedure'' could be applied to a relation in \(\mathcal{C}\) and its negation; and, as we shall see, the possibility of resolving the question in this way is excluded, so far as the usual mathematical theories are concerned (*).

It is in fact in the question of consistency of mathematical theories --- the origin and the very heart of mathematics --- that the results have turned out to be most deceptive. During the years 1920-1930, Hilbert and his school developed new methods for attacking these problems. Having proved the consistency of some partial formalisms, covering a part of arithmetic (cf.~{[}42{]}, {[}43 c{]}), they thought they were on the verge of victory and that they would succeed in proving the consistency not only of arithmetic, but also of set theory, when Gödel used the non-categoricity of arithmetic to deduce that it is impossible to prove, by means of Hilbert's ``finite procedures'', the consistency of any theory E which contains arithmetic (†).

Nevertheless, Gödel's theorem does not close the door altogether to attempts to prove consistency, provided that Hilbert's restrictions concerning ``finite procedures'' are (at least partially) abandoned. Thus, in 1936, Gentzen {[}47{]} succeeded in proving the consistency of formalized arithmetic by ``intuitive'' use of transfinite induction as far as the countable ordinal \(\varepsilon_0\) (Chapter III, § 6, Exercise 14) (**). The degree of ``certainty'' which can be attributed to such an argument is undoubtedly less than for those which satisfy Hilbert's original restrictions, and is essentially a matter depending on the personal psychology of the individual mathematician. But it is no less true that similar ``proofs'' using ``intuitive'' transfinite induction as far as some given ordinal would be considered as an important step forward if they could be applied, for example, to the theory of real numbers or to an important part of set theory.

Furthermore, within set theory itself there are many problems of ``relative'' consistency connected with the many ``hypotheses'' at large in the theory. The most remarkable result in this area is again due to Gödel. In 1940 he proved that, if the theory whose axioms are those of the von Neumann-Bernays system, except for the axiom of choice, is consistent, then the theory obtained by adjoining to these axioms a very strong form of the axiom of choice (of the type which the use of the symbol \(\tau\) gives us) and the generalized continuum hypothesis is also consistent {[}44 b{]}. More recently, I. Novak-Gal has shown that the consistency of the Zermelo-Fraenkel system implies that of the von Neumann-Bernays-Gödel system (*).

